\documentclass[11pt]{article}
% ===========================================================================
%  Shared preamble for all MPM2D worksheets — input right after \documentclass.
%  (Files beginning with "_" are skipped by mpm2d-worksheet-publish.mjs.)
% ===========================================================================
\usepackage[letterpaper,top=2.2cm,bottom=1.9cm,left=1.6cm,right=1.6cm,headheight=28pt,footskip=24pt]{geometry}
\usepackage{amsmath,amssymb}
\usepackage[T1]{fontenc}
\usepackage[utf8]{inputenc}
\usepackage{lmodern}
\usepackage{xcolor}
\usepackage[most]{tcolorbox}
\usepackage{enumitem}
\usepackage{multicol}
\usepackage{fancyhdr}
\usepackage{tikz}
\usetikzlibrary{calc}
\usepackage{pgfplots}
\pgfplotsset{compat=1.18}

\definecolor{exblue}{HTML}{4A90E2}
\definecolor{exbg}{HTML}{E6F3FF}
\definecolor{qorange}{HTML}{E69138}
\definecolor{qbg}{HTML}{FFF7CC}
\definecolor{keygray}{HTML}{475569}
\definecolor{keybg}{HTML}{F1F5F9}
\definecolor{iagreen}{HTML}{14653B}
\definecolor{titleblue}{HTML}{1D4ED8}
% Rotating palette so each worked example gets its own colour.
\definecolor{ex0}{HTML}{2563A0}\definecolor{ex1}{HTML}{3B7D3B}\definecolor{ex2}{HTML}{8A5A00}
\definecolor{ex3}{HTML}{A3327A}\definecolor{ex4}{HTML}{6D28A3}\definecolor{ex5}{HTML}{B08900}
\definecolor{ex6}{HTML}{0E7490}\definecolor{ex7}{HTML}{9A3412}\definecolor{ex8}{HTML}{15803D}

\pagestyle{fancy}
\fancyhf{}
\fancyhead[L]{\footnotesize NAME: \underline{\hspace{3.4cm}}}
\fancyhead[C]{\textbf{Integration Academy}\\[-2pt]{\scriptsize Math Simplified \textperiodcentered{} Success Amplified}}
\fancyhead[R]{\footnotesize MPM2D}
\fancyfoot[L]{\scriptsize dr.merie@IntegrationAcademy.ca}
\fancyfoot[C]{\scriptsize\thepage}
\fancyfoot[R]{\scriptsize IntegrationAcademy.ca}
\renewcommand{\headrulewidth}{1.2pt}
\renewcommand{\footrulewidth}{0.4pt}
\setlength{\parindent}{0pt}
\setlength{\parskip}{4pt}

% Examples are NOT breakable, so each example + its solution stays on one page.
\newtcolorbox{example}[2]{colframe=#1,colback=#1!7,coltitle=white,colbacktitle=#1,fonttitle=\bfseries,title={#2},arc=3pt,boxrule=1pt}
\newtcolorbox{qbox}[1]{colframe=qorange,colback=qbg,coltitle=white,colbacktitle=qorange,fonttitle=\bfseries,title={#1},arc=3pt,breakable,boxrule=1pt}
\newtcolorbox{keybox}{colframe=keygray,colback=keybg,coltitle=white,colbacktitle=keygray,fonttitle=\bfseries,title={$\checkmark$\ Answer Key},arc=3pt,breakable,boxrule=1pt}
\newtcolorbox{recapbox}{colframe=keygray,colback=keybg,coltitle=white,colbacktitle=keygray,fonttitle=\bfseries,title={Question Recap},arc=3pt,breakable,boxrule=1pt}
\newtcolorbox{ideabox}{colframe=titleblue,colback=titleblue!6,coltitle=white,colbacktitle=titleblue,fonttitle=\bfseries,title={The Big Ideas},arc=3pt,breakable,boxrule=1.2pt}
\newtcolorbox{learnbox}{colframe=iagreen,colback=iagreen!5,coltitle=white,colbacktitle=iagreen,fonttitle=\bfseries,title={Concept Essentials},arc=3pt,breakable,boxrule=1.2pt}
\newcommand{\lh}[1]{\par\smallskip{\bfseries\color{iagreen}#1}\par\nobreak\smallskip}

\newcommand{\soln}{\par\smallskip\textit{\textbf{Solution.}}\ }
\newcommand{\workspace}[1]{\par\vspace{3pt}\fcolorbox{black!30}{white}{\begin{minipage}[t][#1][t]{\dimexpr\linewidth-2\fboxsep\relax}\ \end{minipage}}\par}
\newcommand{\sech}[1]{\par\vspace{8pt}{\Large\bfseries\color{iagreen}#1}\par\nobreak\vspace{2pt}{\color{black!20}\hrule height1pt}\vspace{6pt}}

% A compact coordinate plot sized for an example box: \eplot{xmin}{xmax}{ymin}{ymax}{body}
\newcommand{\eplot}[5]{\begin{center}\begin{tikzpicture}
\begin{axis}[width=6.8cm,height=4.4cm,axis lines=middle,xlabel={\tiny$x$},ylabel={\tiny$y$},
grid=both,grid style={gray!16},xmin=#1,xmax=#2,ymin=#3,ymax=#4,
xtick distance=2,ytick distance=2,tick label style={font=\tiny},axis line style={-},samples=120]
#5
\end{axis}\end{tikzpicture}\end{center}}

% A sine-curve plot (degrees on x, 0..360): \splot{ymin}{ymax}{body}
\newcommand{\splot}[3]{\begin{center}\begin{tikzpicture}
\begin{axis}[width=8.6cm,height=4.6cm,axis lines=middle,xlabel={\tiny$x\,(^\circ)$},ylabel={\tiny$y$},
grid=both,grid style={gray!16},xmin=0,xmax=360,ymin=#1,ymax=#2,
xtick distance=90,ytick distance=1,tick label style={font=\tiny},axis line style={-},samples=180]
#3
\end{axis}\end{tikzpicture}\end{center}}

% A small coordinate plot: \plot{xmin}{xmax}{ymin}{ymax}{body}
\newcommand{\plot}[5]{\begin{center}\begin{tikzpicture}
\begin{axis}[width=8.4cm,height=6cm,axis lines=middle,xlabel={\scriptsize$x$},ylabel={\scriptsize$y$},
grid=both,grid style={gray!16},xmin=#1,xmax=#2,ymin=#3,ymax=#4,
xtick distance=2,ytick distance=2,tick label style={font=\tiny},axis line style={-}]
#5
\end{axis}\end{tikzpicture}\end{center}}

% Right triangle (angle theta at bottom-left, right angle at bottom-right):
%   \rtri{drawBase}{drawHeight}{oppLabel}{adjLabel}{hypLabel}
\newcommand{\rtri}[5]{\begin{center}\begin{tikzpicture}[scale=0.85]
\coordinate (A) at (0,0);\coordinate (B) at (#1,0);\coordinate (C) at (#1,#2);
\draw[thick,exblue] (A)--(B)--(C)--cycle;
\draw[black] ($(B)+(-0.32,0)$)--($(B)+(-0.32,0.32)$)--($(B)+(0,0.32)$);
\node[below] at ($(A)!0.5!(B)$){#4};
\node[right] at ($(B)!0.5!(C)$){#3};
\node[above left] at ($(A)!0.5!(C)$){#5};
\node at (0.7,0.22){$\theta$};
\end{tikzpicture}\end{center}}

% General (oblique) triangle with vertex labels and opposite-side labels:
%   \gtri{Alabel}{Blabel}{Clabel}{aSide}{bSide}{cSide}
\newcommand{\gtri}[6]{\begin{center}\begin{tikzpicture}[scale=0.85]
\coordinate (A) at (0,0);\coordinate (B) at (5,0);\coordinate (C) at (1.6,3);
\draw[thick,exblue] (A)--(B)--(C)--cycle;
\node[below left] at (A){#1};\node[below right] at (B){#2};\node[above] at (C){#3};
\node[below] at ($(A)!0.5!(B)$){#6};
\node[right] at ($(B)!0.5!(C)$){#4};
\node[left] at ($(A)!0.5!(C)$){#5};
\end{tikzpicture}\end{center}}

\fancyhead[R]{\footnotesize MCR3U \textperiodcentered{} 5.4}
\begin{document}

\begin{center}
{\LARGE\bfseries\color{titleblue}Worksheet 5.4 --- Sinusoidal Applications}\\[2pt]
{\small Grade 11 Functions, University (MCR3U) \textperiodcentered{} Unit 5: Sinusoidal Functions}
\end{center}

\begin{ideabox}
Model Ferris wheels, tides, and temperature with h = a sin(k(t − d)) + c.
\begin{itemize}[leftmargin=*,itemsep=2pt]
  \item Amplitude $=\dfrac{\text{max}-\text{min}}{2}$; midline $=\dfrac{\text{max}+\text{min}}{2}$.
  \item $k=\dfrac{360^\circ}{\text{period}}$.
  \item Max $=$ midline $+$ amplitude; min $=$ midline $-$ amplitude.
\end{itemize}
\end{ideabox}

\begin{learnbox}
\lh{Modelling real cycles}Turn repeating data (tides, temperatures, wheels) into $y=a\sin\!\left(k(x-d)\right)+c$ using $c=\dfrac{\max+\min}{2}$ and $a=\dfrac{\max-\min}{2}$.
\lh{Setting the period}Find $k$ from the real cycle length: $k=\dfrac{360^\circ}{\text{period}}$, with time in the units given.
\lh{Sine or cosine}Use $-\cos$ to start at a minimum and $+\cos$ to start at a maximum; either models the same motion with a suitable shift.
\end{learnbox}

\sech{Worked Examples}

\begin{example}{ex0}{Example 1 --- Amplitude}
A wheel ranges $2$ m to $12$ m. Amplitude?\soln $\dfrac{12-2}{2}=5$ m.
\end{example}

\begin{example}{ex1}{Example 2 --- Midline}
Same wheel. Midline?\soln $\dfrac{12+2}{2}=7$ m.
\end{example}

\begin{example}{ex2}{Example 3 --- Period to k}
A tide cycle is $12$ h. Find $k$.\soln $k=\dfrac{360^\circ}{12}=30^\circ$/h.
\end{example}

\begin{example}{ex3}{Example 4 --- Build a model}
Temperature ranges $10^\circ$ to $24^\circ$. Amplitude and midline?\soln amplitude $\dfrac{24-10}{2}=7$; midline $\dfrac{24+10}{2}=17$.
\end{example}

\begin{example}{ex4}{Example 5 --- Height at a time}
For $h=5\sin(30t)+7$, find $h$ at $t=3$ (so $30t=90^\circ$).\soln $5\sin90^\circ+7=5(1)+7=12$ m. The model oscillates between $2$ and $12$ about the midline $7$:\splot{0}{14}{\addplot[qorange,dashed,domain=0:360]{7};\addplot[exblue,very thick,domain=0:360]{5*sin(x)+7};}
\end{example}

\begin{example}{ex5}{Example 6 --- Maximum}
Max of $h=5\sin(30t)+7$?\soln midline $+$ amplitude $=7+5=12$ m.
\end{example}

\begin{example}{ex6}{Example 7 --- Minimum}
Min of $h=5\sin(30t)+7$?\soln midline $-$ amplitude $=7-5=2$ m.
\end{example}

\begin{example}{ex7}{Example 8 --- Amplitude}
A wheel ranges $1$ m to $15$ m. Amplitude?\soln $\dfrac{15-1}{2}=7$ m.
\end{example}

\begin{example}{ex8}{Example 9 --- Cycle to k}
A cycle lasts $10$ s. Find $k$.\soln $k=\dfrac{360^\circ}{10}=36^\circ$/s.
\end{example}

\clearpage
\sech{Your Turn --- Show Your Work}
For each question, show all steps clearly in the space provided.

\begin{qbox}{Question 1}
A wheel ranges $2$ m to $18$ m. Find the amplitude.\workspace{2.4cm}
\end{qbox}

\begin{qbox}{Question 2}
Same wheel. Find the midline.\workspace{2.4cm}
\end{qbox}

\begin{qbox}{Question 3}
A cycle lasts $8$ s. Find $k$.\workspace{2.4cm}
\end{qbox}

\begin{qbox}{Question 4}
For $h=6\sin(30t)+8$, find the maximum height.\workspace{2.4cm}
\end{qbox}

\begin{qbox}{Question 5}
For $h=6\sin(30t)+8$, find the minimum height.\workspace{2.4cm}
\end{qbox}

\begin{qbox}{Question 6}
Temperature ranges $8^\circ$ to $22^\circ$. Find the amplitude.\workspace{2.4cm}
\end{qbox}

\begin{qbox}{Question 7}
Same temperature. Find the midline.\workspace{2.4cm}
\end{qbox}

\begin{qbox}{Question 8}
A tide has period $12$ h. Find $k$.\workspace{2.4cm}
\end{qbox}

\begin{qbox}{Question 9}
For $h=5\sin(30t)+7$, find $h$ at $t=3$.\workspace{2.6cm}
\end{qbox}

\begin{qbox}{Question 10}
A wheel ranges $1$ m to $15$ m. Find the amplitude.\workspace{2.4cm}
\end{qbox}

\begin{qbox}{Question 11}
State the range of $h=4\sin(30t)+10$.\workspace{2.4cm}
\end{qbox}

\begin{qbox}{Question 12}
If $k=30^\circ$/h, what is the period?\workspace{2.4cm}
\end{qbox}

\begin{qbox}{Question 13 --- Challenge}
A Ferris wheel turns once in $40$ s; height ranges $1$ m to $21$ m. Find the amplitude, midline, and $k$.\workspace{4cm}
\end{qbox}

\clearpage
\sech{Question Recap \& Answer Key}

\begin{recapbox}
\begin{multicols}{2}
\small
\begin{enumerate}[leftmargin=*,itemsep=3pt]
  \item A wheel ranges $2$ m to $18$ m. Find the amplitude.
  \item Same wheel. Find the midline.
  \item A cycle lasts $8$ s. Find $k$.
  \item For $h=6\sin(30t)+8$, find the maximum height.
  \item For $h=6\sin(30t)+8$, find the minimum height.
  \item Temperature ranges $8^\circ$ to $22^\circ$. Find the amplitude.
  \item Same temperature. Find the midline.
  \item A tide has period $12$ h. Find $k$.
  \item For $h=5\sin(30t)+7$, find $h$ at $t=3$.
  \item A wheel ranges $1$ m to $15$ m. Find the amplitude.
  \item State the range of $h=4\sin(30t)+10$.
  \item If $k=30^\circ$/h, what is the period?
  \item A Ferris wheel turns once in $40$ s; height ranges $1$ m to $21$ m. Find the amplitude, midline, and $k$.
\end{enumerate}
\end{multicols}
\end{recapbox}

\begin{keybox}
\begin{multicols}{2}
\small
\begin{enumerate}[leftmargin=*,itemsep=3pt]
  \item $8$
  \item $10$
  \item $45^\circ$/s
  \item $14$ m
  \item $2$ m
  \item $7$
  \item $15$
  \item $30^\circ$/h
  \item $12$ m
  \item $7$
  \item $6$ to $14$
  \item $12$ h
  \item amplitude $10$, midline $11$, $k=9^\circ$/s
\end{enumerate}
\end{multicols}
\end{keybox}

\end{document}
